\NativeTitle{69}{第八单元素养评价（B）}{（图形的运动）}{40}{6}{一 & 二 & 三 & 四 & 五 & 六}{ &  &  &  &  & }
\NativeSlot{277}{56}{%
\ExamGroup{一、填空。{\normalfont\kaishu\fontsize{10.8bp}{14bp}\selectfont （每空1分，共23分）}}

% p069_s01

\ExamQuestion{1.}{5}%
将一个图形沿一条\mbox{（\hspace{1.8em}）}对折，如果两侧的部分能够完全重合，这样的图形叫作\mbox{（\hspace{1.8em}）}，折痕所在的直线叫作它的\mbox{（\hspace{1.8em}）}。对应点到对称轴的距离\mbox{（\hspace{1.8em}）}，对应点的连线与对称轴互相\mbox{（\hspace{1.8em}）}。\par
}
\NativeSlot{377}{36}{%
% p069_s02

\ExamQuestion{2.}{5}%
图形平移时，图形的\mbox{（\hspace{1.8em}）}和\mbox{（\hspace{1.8em}）}不发生改变，只改变\mbox{（\hspace{1.8em}）}。\par
}
\NativeSlot{402}{95}{%
% p069_s03

\ExamQuestion{3.}{5}%
下图中，是平移现象画“△”，是旋转现象画“○”。图示依次是：滑动解锁、坐摩天轮、开冰箱门、坐缆车。\par
\BookFigure{figures/p069_s03-01.tex}
}
\NativeSlot{567}{56}{%
% p069_s04

\ExamQuestion{4.}{5}%
如图，李阿姨将2千克物品放在台秤上，指针从0刻度开始顺时针旋转\(90^\circ\)，再放2千克物品，指针会\mbox{（\hspace{1.8em}）}时针再旋转\mbox{（\hspace{1.8em}）}度；如果李阿姨把4千克物品全部拿走，那么指针会\mbox{（\hspace{1.8em}）}时针旋转\mbox{（\hspace{1.8em}）}度。\par
\BookFigure{figures/p069_s04-01.tex}
}
\NativeSlot{668}{65}{%
% p069_s05

\ExamQuestion{5.}{5}%
如图，若AB的长度是4厘米，则A点到对称轴的距离是\mbox{（\hspace{1.8em}）}厘米，点C和点F到对称轴的距离都是\mbox{（\hspace{1.8em}）}厘米，点D的对称点是点\mbox{（\hspace{1.8em}）}。\par
\BookFigure{figures/p069_s05-01.tex}
}
